\section{Hyperbolic geometry of the angular pair and recovery of action}
\label{rec:radial}

\subsection{Radion, action ellipse, and marked charts}

The oriented section of phase, action, and memory termed the \emph{radion} jointly preserves the angular increment \(\theta=1\), the signed action \(\sigma\hbar_{\rm ret}\), membership in the revolution \(2\pi_{\HMT}\), the sheet, the nonadic return, and the refinement cylinder. Its unit gluing is expressed in the continuum chart by \(1=S_E-\Phi_T+\varepsilon_1\). The projection retaining only \(\theta=1\) produces the ordinary radian; the enriched section also preserves the register needed to transport and invert its other readings. In this section, \emph{radius} denotes a positive shape coordinate, distinct from that enriched object.

To specify the terms of the gluing identity, the dodecaphase wheel has centres \(30m-10\) in its degree chart. The second phase contributes \(50\), and the representation \(A=1000\alpha\) determines
\[
 S_E=\frac{2\pi}{360}(50+1000\alpha).
\]
The active sector of the tritic selector comprises \(H_+=\{1,4,7\}\) and \(H_-=\{2,5,8\}\). The fifteen unordered pairs, read at the six active positions, produce \(N_6=6\binom62=90\); their extension to eight positions produces \(N_8=8\binom62=120\). The two sheets of the active sector, normalized by the ambient fibre of \(729\) elements, determine
\[
 \rho_{\rm tw}=\frac{2N_6}{729}=\frac{20}{81},\qquad
 \Phi_T=\frac{20}{81}\Delta_4,\qquad
 S_T=1+\Phi_T,\qquad \varepsilon_1=S_T-S_E,
\]
where \(\Delta_4\) is the torsional coordinate in \eqref{rec:eq:electron-basal}. The subtraction retains the preceding sections and directly proves
\(S_E-\Phi_T+\varepsilon_1=S_E-\Phi_T+1+\Phi_T-S_E=1\).
The coefficient comes from the census and the two sheets, not from the value of the final difference.

The angular representations of the preceding sections satisfy, in the chart under consideration,
\[
 A>0,\qquad |C^*|<A,\qquad
 \xi=\frac{C^*}{A},\qquad \eta_{\rm el}=\operatorname{artanh}\xi.
\]
The notation \(\eta_{\rm el}\) distinguishes elliptic anisotropy from the return coefficient \(\eta_{\rm ret}\). For a positive action section \(\hbar\), the semiaxes are
\begin{equation}
 a_S=\sqrt{2\hbar}\,e^{\eta_{\rm el}/2},\qquad
 b_S=\sqrt{2\hbar}\,e^{-\eta_{\rm el}/2},\qquad a_Sb_S=2\hbar.
 \label{rec:eq:elipse}
\end{equation}
The product fixes the action per revolution, whereas the ratio fixes the shape. Define
\[
 R_\star=\sqrt{b_S/a_S}=e^{-\eta_{\rm el}/2},\qquad
 \tau_{\rm el}=iR_\star^2.
\]

\begin{theorem}[Compatibility of the radial, Poincar\'e, and Klein charts]
\label{rec:thm:klein}
Let
\[
 w=\frac{\tau_{\rm el}-i}{\tau_{\rm el}+i},\qquad
 x_K=\frac{2w}{1+w^2}.
\]
On the real diameter of the Poincar\'e disk and its image in the Klein disk,
\begin{align}
 w&=\frac{R_\star^2-1}{R_\star^2+1}
      =-\tanh\frac{\eta_{\rm el}}2,\label{rec:eq:poincare}\\
 x_K&=\frac{R_\star^4-1}{R_\star^4+1}
      =-\tanh\eta_{\rm el}=-\frac{C^*}{A}.\label{rec:eq:klein}
\end{align}
All these transformations are invertible on their domains; in particular,
\(R_\star^2=(1+w)/(1-w)\).
\end{theorem}
\begin{proof}
Substitute \(\tau=iR_\star^2\) into the fractional transformation. The simplified ratio has denominator \(R_\star^2+1>0\) and \(-1<w<1\). Substituting this ratio into \(2w/(1+w^2)\) and multiplying numerator and denominator by \((R_\star^2+1)^2\) gives \((R_\star^4-1)/(R_\star^4+1)\). The identity \(R_\star^4=e^{-2\eta_{\rm el}}\) yields the hyperbolic tangent. Finally, \(\tanh\eta_{\rm el}=\xi\). The recovery maps follow by solving the ratios.
\end{proof}

The sign in \eqref{rec:eq:klein} corresponds to the chosen orientation of the diameter. The opposite orientation changes both signs. This is a hyperbolic coordinate of the marked rectangular chart: it retains the marking and is not confused with a modular form or with a coordinate on the modular quotient stripped of its orientation.

Exchanging the semiaxes acts by
\begin{equation}
 (R_\star,\eta_{\rm el},\tau_{\rm el},w,x_K)
 \longmapsto(R_\star^{-1},-\eta_{\rm el},-1/\tau_{\rm el},-w,-x_K).
 \label{rec:eq:reciproco-radial}
\end{equation}
On this diameter, the hyperbolic metrics of curvature \(-1\) satisfy
\begin{equation}
 \frac{|d\tau|^2}{(\Im\tau)^2}
 =\frac{4\,dw^2}{(1-w^2)^2}
 =\frac{dx_K^2}{(1-x_K^2)^2}
 =d\eta_{\rm el}^2=4\frac{dR_\star^2}{R_\star^2}.
 \label{rec:eq:metricas}
\end{equation}
These identities are verified by differentiating \eqref{rec:eq:poincare}--\eqref{rec:eq:klein}. In particular, \(d(0,w)=|\eta_{\rm el}|\) and \(d(w,-w)=2|\eta_{\rm el}|\). The transformation \(w\mapsto-w\) is a rotation by \(\pi\) on the entire disk and reflection about the centre on the generated diameter.

The normalization \((Q/a_S,P/b_S)\) of the ellipse's interior produces another realization of the Klein disk, with position coordinates. Its metric is
\[
 ds_K^2=\frac{(1-r^2)(du^2+dv^2)+(u\,du+v\,dv)^2}{(1-r^2)^2},
 \qquad r^2=u^2+v^2<1.
\]
This physical chart and the marked shape chart retain distinct domains. On the focal branch \(\eta_{\rm el}\ge0\), the eccentricity \(e_f=\sqrt{1-e^{-2\eta_{\rm el}}}\) gives depth \(u_f=\operatorname{artanh}e_f\), whence \(\eta_{\rm el}=\log\cosh u_f\). Focal depth is not identified with anisotropy.

\begin{figure}[htbp]
\centering
\begin{tikzpicture}[>=Latex,scale=1.0,every node/.style={font=\small}]
\draw[ink,thick] (0,0) circle (1.55);
\draw[->] (-1.8,0)--(1.85,0);
\fill[accent] (-.75,0) circle (2pt) node[below=4pt] {\(w\)};
\fill[accent] (.75,0) circle (2pt) node[below=4pt] {\(-w\)};
\draw[<->,accent,thick] (-.75,.30)--(.75,.30);
\node at (0,-2.0) {Shape chart: \(w\mapsto-w\)};
\begin{scope}[xshift=7cm]
\draw[ink,thick] (0,0) circle (1.15);
\draw[->] (-1.6,0)--(2.05,0);
\draw[->] (0,-1.4)--(0,1.6);
\draw[accent,thick] (0,0)--(30:.65);
\draw[accent,thick,dashed] (30:.65)--(30:2.03);
\fill[accent] (30:.65) circle (2pt) node[above=3pt] {\(z\)};
\fill[accent] (30:2.03) circle (2pt) node[above=3pt] {\(1/\overline z\)};
\node at (0,-2.0) {Spectral chart: radial inversion};
\end{scope}
\end{tikzpicture}
\caption{Two typed reciprocities. The first remains inside the disk and transports shape orientation; the second exchanges the interior and exterior of the unit circle. The plotted positions are schematic.}
\label{fig:rec-dos-cartas}
\end{figure}

\subsection{Complete substitution of the angular representations}
\label{rec:sustitucion}

The dodecaphasic closure already constructed expresses the same coordinate \(\alpha\) through
\[
 \alpha=\pi_{\HMT}+e_{\HMT}-\varphi_{\HMT}-4-\kappa_{\rm ag},
 \qquad \kappa_{\rm ag}=\frac{K_{\rm int}}{1000^{12}-1}.
\]
The compatible base, phase origin, and carry boundary are retained. The article's complete analytic recurrence represents that correlated coordinate; its representation is not replaced by a finite jet. Abbreviating the HMT subscripts, the action readers of \eqref{iv:action:h5}--\eqref{iv:action:eta-ret} are
\begin{align}
 H_5&=\frac12\sqrt{\frac{1000\alpha}{\varphi}}-\alpha
       +\frac9{16}\alpha^2-\frac59\alpha^3
       +\frac7{48}\alpha^4-\frac1{54}\alpha^5,\nonumber\\
 D_A&=e^{-100\pi\alpha/9},\qquad
 \mathcal C_\pi=169\alpha^6+\frac{D_A\alpha^7}{1-D_A\alpha},\nonumber\\
 \eta_{\rm ret}&=H_5-\frac{90}{\pi}\mathcal C_\pi,
 \qquad A=1000\alpha,\qquad C^*=2(\eta_{\rm ret}+\alpha).
 \label{rec:eq:lectores-completos}
\end{align}
The rational continuation retains every term of the regional recurrence. Its convergence follows without truncation: for \(\alpha>0\),
\[
 0<\alpha e^{-100\pi\alpha/9}\le\frac9{100\pi e}<1,
\]
since \(xe^{-cx}\) attains its maximum \(1/(ce)\) at \(x=1/c\).

\begin{proposition}[Structural cancellation in the angular contrast]
\label{rec:prop:sustitucion}
The generated ratio \(\xi=C^*/A\) satisfies exactly
\begin{align}
 \xi={}&\frac1{\sqrt{1000\varphi\alpha}}
 +\frac{9\alpha}{8000}-\frac{\alpha^2}{900}
 +\frac{7\alpha^3}{24000}-\frac{\alpha^4}{27000}\nonumber\\
 &-\frac9{50\pi}
 \left(169\alpha^5+\frac{D_A\alpha^6}{1-D_A\alpha}\right).
 \label{rec:eq:xi-completa}
\end{align}
\end{proposition}
\begin{proof}
Start from \(\xi=(\eta_{\rm ret}+\alpha)/(500\alpha)\). The term \(-\alpha\) in \(H_5\) cancels against the \(+\alpha\) of the counter-angle. Dividing each remaining term by \(500\alpha\) yields the expression: the square root gives \(1/\sqrt{1000\varphi\alpha}\), and \(90/500=9/50\). The regional continuation retains its complete denominator.
\end{proof}

Consequently, \(x_K=-\xi\) and \(R_\star=((1-\xi)/(1+\xi))^{1/4}\) are explicit expressions in the regional representations and the register. Substituting the closure of \(\alpha\) completes the dependence on \(\pi,\varphi,e,\kappa_{\rm ag}\), without promoting metrological identifications to generators.

\subsection{Recovery of Planck's action with orientation preserved}

\begin{theorem}[Radial recovery map for the return section]
\label{rec:thm:accion}
In the positive chart, if \(q=R_\star^4\), then
\begin{align}
 \eta_{\rm ret}&=\alpha\frac{499-501q}{1+q},\label{rec:eq:eta-radio}\\
 \hbar_{\rm ret}&=\alpha\frac{499-501q}{1+q}\,10^{-34}\mathcal U_S,
 \qquad\hbar_{\rm pre}=H_5\,10^{-34}\mathcal U_S,\label{rec:eq:planck-radio}\\
 R_{\rm act}&=\frac{\hbar_{\rm pre}}{\hbar_{\rm ret}}
 =\frac{H_5(1+q)}{\alpha(499-501q)}.\label{rec:eq:ratio-radio}
\end{align}
The branch \(\eta_{\rm ret}>0\) is equivalent to \(q<499/501\), within \(q>0\).
\end{theorem}
\begin{proof}
The angular identities give \(\eta_{\rm ret}=\alpha(500\xi-1)\). Substitute \(\xi=(1-q)/(1+q)\); the numerator is \(499-501q\). Both the denominator and \(\alpha\) are positive. The dimensional realization and the ratio are those of the action sections already constructed. The decade \(-34\) and the unit \(\mathcal U_S\) retain their preceding determinantal derivation.
\end{proof}

This formula makes the Planck section within the counter-angle explicit. In terms of the action quantities, the angular relation is
\[
 C^*=2\left(\frac{\hbar_{\rm ret}}{10^{-34}\mathcal U_S}+\alpha\right).
\]
Division by the dimensional basis extracts the dimensionless coefficient; it does not add an action quantity to a dimensionless constant.

For the angular orientation \(\varepsilon\in\{-1,+1\}\), define
\[
 C^*_{\varepsilon}=2\varepsilon(\eta_{\rm ret}+\alpha),\quad
 \xi_\varepsilon=C^*_{\varepsilon}/A,\quad
 q_\varepsilon=\frac{1-\xi_\varepsilon}{1+\xi_\varepsilon}>0.
\]
The recovery map on both charts is
\begin{equation}
 \eta_{\rm ret}=\alpha\left(500\varepsilon
        \frac{1-q_\varepsilon}{1+q_\varepsilon}-1\right).
 \label{rec:eq:eta-orientada}
\end{equation}

\begin{corollary}[Invariance of the conjugate recovery map]
The transformation \((\varepsilon,q)\mapsto(-\varepsilon,q^{-1})\) preserves \(\eta_{\rm ret}\), \(\hbar_{\rm ret}\), and \(R_{\rm act}\).
\end{corollary}
\begin{proof}
The ratio \((1-q)/(1+q)\) changes sign when \(q\) is inverted. Its product with \(\varepsilon\) remains invariant. Apply \eqref{rec:eq:eta-orientada}, followed by the action representations.
\end{proof}

If orientation is discarded and \(\varepsilon=+1\) is incorrectly held fixed, the same expression gives \(\eta_{\rm ret}(q^{-1})=-\eta_{\rm ret}(q)-2\alpha\). The register therefore distinguishes two algebraically different operations. Nor is the angular orientation \(\varepsilon\) automatically identified with the symplectic sign \(\sigma\): exchanging the axes is antisymplectic, whereas the rotation that exchanges them preserves the area form.

\subsection{Intertwining phase and anisotropy}

Let
\[
 S_\eta=\operatorname{diag}(e^{\eta/2},e^{-\eta/2}),\quad
 J_0=\begin{pmatrix}0&-1\\1&0\end{pmatrix},\quad
 J_\eta=S_\eta J_0S_\eta^{-1}
       =\begin{pmatrix}0&-e^\eta\\e^{-\eta}&0\end{pmatrix}.
\]
Then \(\det S_\eta=1\), \(J_\eta^2=-I\), and
\(C_\eta(\theta)=e^{\theta J_\eta}=S_\eta e^{\theta J_0}S_\eta^{-1}\).
The law \(S_{\eta+\delta}=S_\delta S_\eta\) proves the exact intertwining identity
\begin{equation}
 C_{\eta+\delta}(\theta)S_\delta=S_\delta C_\eta(\theta).
 \label{rec:eq:entrelazamiento}
\end{equation}
The transformation preserves \(dQ\wedge dP\) and carries the phase evolution to the new ellipse. With the normalization of \eqref{rec:eq:elipse}, the swept action is \(\hbar\theta\), or \(\sigma\hbar\theta\) on the oriented sheet. The phase \(\theta\), the anisotropy \(\eta_{\rm el}\), and the parameter \(t\) of the flow of \(K\) thus retain distinct functions, related by explicit operators.

\subsection{Radial geodesic and differential realization of the direction of K}

Each radius in \eqref{rec:eq:flujo} determines a rectangular chart and a hyperbolic coordinate:
\[
 \tau_i(t)=ir_i(t)^2,\qquad
 w_i(t)=\frac{r_i(t)^2-1}{r_i(t)^2+1}=\tanh(th_i).
\]

\begin{theorem}[Geodesic of compensated radii]
\label{rec:thm:geodesica}
In the product of twelve hyperbolic diameters,
\begin{equation}
 d(w(s),w(t))=2|t-s|\,|\log\rho|,\qquad
 \sum_i\operatorname{artanh}w_i(t)=0.
 \label{rec:eq:geodesica}
\end{equation}
Composition of advances is given by
\[
 w_i(t+s)=\frac{w_i(t)+w_i(s)}{1+w_i(t)w_i(s)}.
\]
\end{theorem}
\begin{proof}
In coordinates \(a_i=\log r_i\), the product metric is \(4\sum_i da_i^2\). The curve \(a(t)=th\) is a straight line, with speed \(2\|h\|=2|\log\rho|\), and \(\sum_i a_i=0\). Each diameter is totally geodesic. Distance in the product is the square root of the sum of squared distances in its factors, which are \(2|t-s||h_i|\). The vanishing-sum identity and the composition law follow from \(\operatorname{artanh}w_i=th_i\) and the addition formula for the hyperbolic tangent.
\end{proof}

The regional ratio determines the speed of this realization; \(K\) determines its direction. If \(\rho=1\), the trajectory is constant. For \(\rho\ne1\), the space of radii with product one has dimension eleven, with tangent space \(V_{11}=\mathbf1^\perp\), and the normal to the trajectory within it is
\[
 V_{10}(K)=V_{11}\cap u^\perp,\qquad
 P_{10}(K)=P_{11}-\frac{uu^{\mathsf T}}{\|u\|^2}.
\]
The projector of Proposition~\ref{prop:k-reduccion-diez} therefore acquires a differential realization: it removes the uniform scale and distinguishes the generated direction from its ten transverse variations. Its domain is this radial space, not the real twenty-four-dimensional space of all the disks.

With \(D_t=\operatorname{diag}(\operatorname{sech}^2(th_i))\), its expression in coordinates \(w\) is \(D_tP_{10}D_t^{-1}\). Indeed, \(D_t\) is the differential of \(a\mapsto\tanh a\), and transforms the metric \(4I\) into \(4\operatorname{diag}((1-w_i^2)^{-2})\). The transported projection is orthogonal for that metric.

\subsection{Lattice lift and rectangular lift}

The radii admit two realizations that should be retained explicitly. The operator \(G_t=\bigoplus_i r_iI_2\) scales both axes of each \(A_2\) plane simultaneously. It preserves the total covolume and commutes with the ternary isometry. The homothety of each plane preserves its internal complex modulus.

The chart \(\tau_i=ir_i^2\) has another realization, of fixed area:
\[
 T_i=\operatorname{diag}(r_i^{-1},r_i),\qquad
 (\omega_1,\omega_2)=(r_i^{-1},ir_i),\qquad
 \omega_2/\omega_1=ir_i^2.
\]
This is the elliptic lift with \(\eta=-2\log r_i\). This convention transforms physical periods. The fractional linear action of the same matrix on \(i\) would give \(i/r_i^2\); to obtain \(ir_i^2\) under that other convention, one uses \(T_i^{-1}\). Distinguishing the actions prevents an inadvertent inversion.

Both realizations are reconstructed from the same radii, but the ternary commutation of the homothety does not transfer to the anisotropic deformation. The state and its readers allow the two geometries to be compared without identifying their operators.
